A graduate course in algebraic topology from the University of Chicago. 

\begin{definition}[Pointed Topological Space]
  A \textbf{pointed topological space} is a topological space $X$ with a distinguished point $x_0 \in X$, denoted $(X, x_0)$. 
\end{definition}

\begin{definition}[Wedge Sum]
  Given spaces $X, Y$ with chosen points $x_0 \in X, y_0 \in Y$, we define the \textbf{wedge sum} as
	\begin{equation}
		X \vee Y \coloneqq \frac{X \sqcup Y}{\sim}, \quad x_0 \sim y_0
	\end{equation}
\end{definition}


\begin{definition}[Cell]
  Let $X$ be a topological space. An \textbf{$n$-cell} is a subspace $e^n \subset X$ homeomorphic to the closed unit ball
  \begin{equation}
    \mathring{D}^n = \{x \in \mathbb{R}^n : \|x\| \leq 1\}.
  \end{equation}
\end{definition}

Therefore, colloquially, $0$-cells are points, $1$-cells are curves, $2$-cells are disks and rectangles, $3$-cells are spheres and boxes, and so on. 

\begin{definition}[Cell Complex]
  A \textbf{CW complex} is a topological space $X$ constructed inductively as follows.
  \begin{enumerate}
    \item \textit{0-Skeleton}. Start with a discrete set $X^0$, whose points are called \textbf{0-cells}.

    \item \textit{$n$-Skeleton}. Given the $(n-1)$-skeleton $X^{n-1}$, construct the \textbf{$n$-skeleton} $X^n$ by attaching a collection of closed $n$-dimensional disks $\{D_\alpha^n\}_{\alpha \in I_n}$ along continuous \textbf{attaching maps}
      \begin{equation}
        \varphi_\alpha: S^{n-1} = \partial D_\alpha^n \to X^{n-1}.
      \end{equation}
      More precisely, $X^n$ is the quotient space
      \begin{equation}
        X^n = \left( X^{n-1} \sqcup \coprod_{\alpha \in I_n} D_\alpha^n \right) \big/ \sim,
      \end{equation}
      where $x \sim \varphi_\alpha(x)$ for all $x \in \partial D_\alpha^n$. The image of $\operatorname{int}(D_\alpha^n)$ in $X^n$ is called an \textbf{$n$-cell} $e_\alpha^n$.

    \item \textit{Closure-Finiteness (C)}. The closure of every cell $e_\alpha^n$ intersects only finitely many cells of $X$.

    \item \textit{Weak Topology (W)}. Define
      \begin{equation}
        X = \bigcup_{n=0}^\infty X^n
      \end{equation}
      with the weak topology determined by the skeleta: a subset $A \subset X$ is closed if and only if $A \cap X^n$ is closed in $X^n$ for every $n \geq 0$.
  \end{enumerate}
  A space constructed in this way is called a \textbf{CW complex}, where C stands for closure-finiteness and W stands for weak topology.
\end{definition}
